\documentclass[a4paper,11pt,fleqn]{scrartcl}%fleqn pour aligner les equations à gauche
\usepackage{../../Styles/Eval}


\newcommand{\chnum}{Chapitre 5 et 18}
\newcommand{\chtitle}{Mouvement des planètes et lunette astronomique}
\newcommand{\dtype}{DS 5}

\begin{document}
\maketitle
\section{Mars due sous l'\oe il de Kepler}
Johannes Kepler (1571-1630) est un astronome allemand connu pour avoir établi les trois lois sui portent son nom et qui permettent notamment de décrire le mouvement des planètes du système solaire.\\
On étudie dans cet exercice la troisième loi de Kepler appliquée à la planète Mars, puis on détermine les caractéristiques d'une lunette astronomique permettant d'observer cette planète.

\subsection{Etude et utilisation des lois de Kepler}
\begin{data}
    Valeur de l'unité astronomique (UA) : \SI{149597870700}{m}
\end{data}
En 1600, Kepler devient l'assistant de l'astronome danois Tycho Brahe. Celui-ci le charge d'étudier la trajectoire de Mars et de déterminer son orbite à l'aide des observations astronomiques et des mesures qu'il a lui-même effectuées. Cette tâche, longue et délicate, conduit Kepler à envisager une trajectoire elliptique pour Mars, il publie ses deux premières lois en 1609, puis la troisième loi en 1618.\\
\begin{flushright}
    D'après \textit{La science moderne : de 1450 à 1800}, de René Taton
\end{flushright}
\begin{figure}[h!]
    \centering
    \includegraphics[width=.5\linewidth]{Ressources/TSPE_DS5_25_Ellipse.png}
    \caption{Trajectoire elliptique de Mars dans le référentiel héliocentrique (échelle non respectée). Les aires $A_1$ et $A_2$ sont balayées pendant des durées égales.}
    \label{fig1}
\end{figure}
Le référentiel d'étude est le référentiel héliocentrique supposé galiléen : son origine est au centre du Soleil et ses axes sont fixés par rapport aux étoiles lointaines.\\
\begin{enumerate}
    \item Énoncer les deux premières loide de Kepler.
    \item En exploitant la deuxième loi et le schéma de la Figure~\ref{fig1}, justifier que le mouvement de Mars dans le référentiel héliocentrique n'est pas uniforme.
\end{enumerate}
Les périodes de révolution $T$ 'en années) et les demi-grands axes $a$ (en unité astronomique) des trajectoires des planètes du système solaire (à l'exception de Mars) sont saisies dans un programme écrit en langage Python. Afin de vérifier la troisième loi de Kepler. Un extrait de ce programme est donné sur la Figure~\ref{fig2}.
\begin{figure}[h!]
    \centering
    \includegraphics[width=.8\linewidth]{Ressources/TSPE_DS5_25_Python.png}
    \caption{Extrait du programme Python permettant de vérifier la troisième loi de Kepler.}
    \label{fig2}
\end{figure}
Le programme permet d'obtenir une représenation graphique dont un zoom est proposéen Figure~\ref{fig3}.
\begin{figure}[h!]
    \centering
    \includegraphics[width=.8\linewidth]{Ressources/TSPE_DS5_25_Graph.png}
    \caption{Modélisation de la troisième loi de Kepler à partir des données des planètes du système solaire.}
    \label{fig3}
\end{figure}

Dans le tableau 1 sont regroupés les périodes de révolution (en s) et les demi-grands axes (en m) des trajectoires des plnètes du système solaire (à l'exception de Mars).
\begin{center}
    \begin{tabular}{|c|c|c|}
        \hline
        Planète & Période de révolution $T$ (en s) & Demi-grand axe $a$ (en m) \\
        \hline
        Mercure & \num{7.60e6} & \num{5.79e10} \\
        Vénus & \num{1.94e7} & \num{1.08e11} \\
        Terre & \num{3.15e7} & \num{1.50e11} \\
        Jupiter & \num{3.74e8} & \num{7.78e11} \\
        Saturne & \num{9.29e8} & \num{1.43e12} \\
        Uranus & \num{2.66e9} & \num{2.87e12} \\
        Neptune & \num{5.20e9} & \num{4.50e12} \\
        \hline
    \end{tabular}
    \captionof{table}{Périodes de révolution et demi-grands axes des trajectoires des planètes du système solaire (à l'exception de Mars).}
\end{center}
\begin{enumerate}[resume]
    \item Recopier sur la copie la ligne 8 du programme et la compléter.
    \item Proposer sur la copie un commentaire en précisant la finalité des lignes 16 et 17 du programme et les unités des grandeurs calculées.
    \item Commenter le Figure~\ref{fig3} au regard des lois de Kepler.
\end{enumerate}

À partir des relevés de Tycho Brahe, Kepler a pu déterminer que la période de révolution de Mars, notée $T_{Mars}$ était de 687 jours.

\begin{enumerate}[resume]
    \item Déterminer la valeur du demi-grand axe de l'orbite de Mars, noté $a_{Mars}$, et justifier qu'elle correspond à la quatrième planète du système solaire en partant du Soleil.
\end{enumerate}
\subsection{Observer Mars à l'aide d'une lunette astronomique}
En 1611, Johannes Jepler propose dans son ourvrage \textit{Dioptricae} une nouvelle combinaison optique pour la lunette astronomique utilisée par Galilée en remplaçant la lentille divergente par une lentille convergente.
\begin{flushright}
    D'après \textit{La science moderne : de 1450 à 1800}, de René Taton
\end{flushright}
Fondées sur ce modèle, les lunettes astronomiques actuelles sont formées de deux lentilles minces convergentes. On a représenté sur la Figure~\ref{fig4}, le schéma optique d'une lunette astronomique afocale.
\begin{data}
    \begin{itemize}
        \item distance minimale entre Mars et la Terre : \SI{62,07e6}{km}
        \item à l'\oe il nu, l'angle sous lequel est vue la Lune est de \SI{9,0e-3}{rad}
        \item le diamètre de Mars est de \SI{6794}{km}
        \item pour des angles suffisamment petits, c'est-à-dire très inférieurs à 1 radian, on peut écrire : $$\tan(\theta) \approx \theta \text{~ où } \theta \text{ est en radian}$$
    \end{itemize}
\end{data}
On présente sur le schéma ci-dessous l'angle $\theta$ sous lequel l'observateur voit un objet $AB$ de hauteur $h$ lorsqu'il se situe à une distance $d$ grande devant $h$.
\begin{figure}[h!]
    \centering
    \includegraphics[width=.8\linewidth]{Ressources/TSPE_DS5_25_Observation.png}
    \label{fig5}
\end{figure}

\begin{enumerate}[resume]
    \item Sur la Figure~\ref{fig4}, on note $L_1$ et $L_2$ les deux lentilles minces convergentes. Préciser la lentille correspondant à l'objectif et celle correspondant à l'occulaire de la lunette.
\end{enumerate}

On considère un objet situé à l'infini, noté $A_{\infty}B_{\infty}$. On observe cet objet avec la lunette.

\begin{enumerate}[resume]
    \item Tracer sur la Figure~\ref{fig4} la marche des rayons lumineux provenent de $B_{\infty}$ à travers la lentille $L_1$ et la lentille $L_2$ en faisant apparaitre l'imahe intermédiaire, notée $A_1B_1$ de l'objet $A_{\infty}B_{\infty}$ formée par la lentille $L_1$.
\end{enumerate}

La plupart du temps, il est difficile d'observer Mars depuis la Terre, notamment à cause de sa petite taille. La situation la plus favorable est quand Mars est en opposition, c'est-à-dire alignée avec la Terre et le Soleil. Cette situation correspond au moment où Mars est au plus près de la Terre.
\begin{flushright}
    D'après \textit{http://www.observatoiredeparis.psl.eu}
\end{flushright}
On observe Mars à l'aide d'une lunette astronomique dont les caractéristiques sont données au tableau 2.
\begin{center}
    \begin{tabular}{|c|c|}
        \hline
        Distance focale de l'objectf & \SI{900}{mm} \\
        \hline
        Diamètre de l'objectif & \SI{70}{mm} \\
        \hline
        Masse du tube optique & \SI{1,75}{kg} \\
        \hline
        Distance focale des oculaires interchangeables & \SI{10}{mm}, \SI{25}{mm} et \SI{40}{mm}\\
        \hline
    \end{tabular}
    \captionof{table}{Fiche technique d'une lunette astronomique 70/900.}
\end{center}

\begin{enumerate}[resume]
    \item Représenter les angles $\theta_{1}$ (angle sous lequel est vu Mars à l'\oe il nu) et $\theta_2$ (angle sous lequel est vu Mars à l'aide de la lunette) sur la Figure~\ref{fig4}.
    \item On rappelle que le grossissement $G$ de la lunette s'écrit $G = \dfrac{\theta_2}{\theta_1}$. Établir que le grossissement s'exprime en fonction des distance focales de l'objectif et de l'oculaire notée respectivement $f'_{obj}$ et $f'_{ocu}$ : $G = \dfrac{f'_{obj}}{f'_{ocu}}$.
    \item Dans la situation où Mars est au plus près de la Terre, déterminer parmi les oculaires fournies avec la lunette décrite dans le tableau 2, celui qui permet à un observateur de voir Mars au moins aussi grosse que la Lune vue à l'\oe il nu.
\end{enumerate}

\begin{figure}[b]
    \centering
    \includegraphics[width=\linewidth]{Ressources/TSPE_DS5_25_Lunette.png}
    \caption{Schéma optique d'une lunette astronomique afocale.}
    \label{fig4}
\end{figure}
\end{document}